\section{A linguagem de conjuntos e lógica básica}
A linguagem de conjuntos é uma ferramenta poderosa para a construção de teorias matemáticas rigorosas modernas.
A lógica de primeira ordem é uma linguagem básica sobre qual a teoria dos conjuntos, e, portanto, o cálculo diferencial e integral, pode ser formalizado.

Conjuntos podem ser pensados como coleções de objetos.
Tais objetos, por sua vez, também são conjuntos.

Utilizamos letras como $x, y, z, a, b, c$, etc, para denotar objetos (conjuntos).
Formalmente, tais símbolos denotam o que chamamos de \emph{variáveis lógicas}.
Intuitivamente, eles podem assumir o valor de qualquer objeto (conjunto) que desejarmos.

Ao longo do desenvolvimento da teoria, símbolos que denotam objetos ``fixos'' serão inseridos, como o conjunto vazio $\emptyset$, o conjunto dos números reais $\mathbb R$, e os números reais $0, 1$ e $\pi$, por exemplo.
Tais símbolos são chamados de \emph{constantes} e são introduzidos por meio de definições.
Adiaremos a discussão sobre o significado de definir algo matematicamente para quando tal fenômeno ocorrer na prática, uma vez que tal discussão será mais natural naquela altura.

Escrevemos $x\in y$ para representar a ideia de que $x$ é um elemento do conjunto $y$.
Lemos $x\in y$ como ``$x$ pertence a $y$'' ou ``$x$ é um elemento de $y$''.

Escrevemos $x=y$ para dizer que $x$ e $y$ são o mesmo objeto.
Lemos $x=y$ como ``$x$ é igual a $y$''.

O princípio lógico da substituição se aplica: se $x\in y$ e $x=z$, por exemplo, então podemos concluir que $z\in y$, uma vez que intuitivamente $x$ e $z$ são o mesmo objeto.

As afirmações básicas sobre conjuntos são expressas por meio de fórmulas lógicas.
As fórmulas mais simples são as do tipo $x\in y$ e do tipo $x=y$.
Para obter fórmulas mais complexas, podemos usar os conectivos lógicos $\neg$ (negação), $\wedge$ (conjunção), $\vee$ (disjunção), $\to$ (condicional) e $\leftrightarrow$ (bicondicional), bem como os quantificadores $\forall$ (quantificador universal) e $\exists$ (quantificador existencial).

No geral, denotamos fórmulas por letras gregas como $\phi$ e $\psi$ (lê-se: phi e psi, respectivamente).

Fórmulas podem ser verdadeiras ou falsas.

\subsection{A negação de uma fórmula}
Se $\phi$ é uma fórmula, então $\neg \phi$ é a negação de $\phi$.
Intuitivamente, $\neg \phi$ é verdadeira se e somente se $\phi$ é falsa.

Se $\phi$ é a fórmula $x\in y$, então $\neg \phi$ é a fórmula $\neg(x\in y)$.
Intuitivamente, se $x$ é um elemento de $y$, então $x\in y$ é verdadeira, e $\neg(x\in y)$ é falsa.
Por outro lado, se $x$ não é um elemento de $y$, então $x\in y$ é falsa, e $\neg(x\in y)$ é verdadeira.

Uma notação mais concisa para escrever $\neg(x\in y)$ é $x\notin y$.
Lêmos $x\notin y$ como ``$x$ não pertence a $y$'' ou ``$x$ não é um elemento de $y$''.

Da mesma forma, uma notação mais concisa para escrever $\neg(x=y)$ é $x\neq y$.
Lêmos $x\neq y$ como ``$x$ é diferente de $y$'' ou ``$x$ não é igual a $y$''.

A tabela abaixo resume a relação entre o valor de verdade de uma fórmula $\phi$ e sua negação $\neg \phi$.
``V'' representa o valor de verdade verdadeiro, e ``F'' representa o valor de verdade falso.

\begin{center}
    \begin{tabular}{c|c}
        $\phi$ & $\neg \phi$ \\
        \hline
        V & F \\
        F & V
    \end{tabular}
\end{center}

\begin{example}
    Para ilustrar o comportamento do conectivo de negação, tomemos um exemplo ``cotidiano''.
    Considere a afirmação ``a Terra é plana''.
    Sua negação, ``$\neg$(a Terra é plana)'', é a afirmação ``a Terra não é plana''.
    Uma vez que a afirmação ``a Terra é plana'' é falsa, sua negação ``a Terra não é plana'' é verdadeira.
\end{example}

\subsection{A conjunção de fórmulas}
Se $\phi$ e $\psi$ são fórmulas, então $\phi\wedge \psi$ é a conjunção de $\phi$ e $\psi$.
Lemos $\phi\wedge \psi$ como ``$\phi$ e $\psi$''.
Também podemos escrever $\phi\wedge \psi$ como literalmente $\phi$ e $\psi$.

\begin{example}
    A conjunção de $x\in y$ com $x=y$ é a fórmula
    \[
        x\in y\wedge x=y.
    \]
    Tal fórmula pode ser escrita também como
    \[
        x\in y \text{ e } x=y.
    \]
\end{example}

\begin{example}
    A conjunção de $x\in y$ com $x=y$ é a fórmula
    \[
        x\in y\wedge x=y.
    \]
    Tal fórmula pode ser escrita também como
    \[
        x\in y \text{ e } x=y.
    \]
\end{example}

Para que a fórmula $\phi\wedge \psi$ seja verdadeira, é necessário que ambas as fórmulas $\phi$ e $\psi$ sejam verdadeiras.
Caso ao menos uma das fórmulas $\phi$ ou $\psi$ seja falsa, então $\phi\wedge \psi$ é falsa.

A tabela abaixo resume a relação entre os valores de verdade de $\phi$, $\psi$ e $\phi\wedge \psi$.

\begin{center}
    \begin{tabular}{cc|c}
        $\phi$ & $\psi$ & $\phi\wedge \psi$ \\
        \hline
        V & V & V \\
        V & F & F \\
        F & V & F \\
        F & F & F
    \end{tabular}
\end{center}

Façamos um exemplo ``cotidiano'' para ilustrar o conceito de conjunção.

\begin{example}
    Considere a afirmação ``existem patos em Vênus'' e a afirmação ``Os direitos e garantias individuais são uma cláusula pétrea da Constituição Federal de 1988''.

    A conjunção dessas afirmações é a afirmação ``existem patos em Vênus e os direitos e garantias individuais são uma cláusula pétrea da Constituição Federal de 1988''.
    Apesar de uma das afirmações que compõe a conjunção ser verdadeira, a conjunção é falsa, uma vez que a outra é falsa.
\end{example}

\subsection{A disjunção de fórmulas}
Se $\phi$ e $\psi$ são fórmulas, então $\phi\vee \psi$ é a disjunção de $\phi$ e $\psi$.
Lêmos $\phi\vee \psi$ como ``$\phi$ ou $\psi$''.
Também podemos escrever $\phi\vee \psi$ como literalmente $\phi$ ou $\psi$

\begin{example}
    A disjunção da fórmula $x\in y\wedge x=y$ com a fórmula $x\in z$ é a fórmula
    \[
        \left(x\in y\wedge x=y\right) \vee (x\in z).
    \]
    Tal fórmula pode ser escrita também como
    \[
        (x\in y \text{ e } x=y) \text{ ou } x\in z.
    \]
\end{example}

Para que a fórmula $\phi\vee \psi$ seja verdadeira, basta que ao menos uma das fórmulas $\phi$ ou $\psi$ seja verdadeira.
Caso ambas as fórmulas $\phi$ e $\psi$ sejam falsas, então $\phi\vee \psi$ é falsa.

A tabela abaixo resume a relação entre os valores de verdade de $\phi$, $\psi$ e $\phi\vee \psi$.

\begin{center}
    \begin{tabular}{cc|c}
        $\phi$ & $\psi$ & $\phi\vee \psi$ \\
        \hline
        V & V & V \\
        V & F & V \\
        F & V & V \\
        F & F & F
    \end{tabular}
\end{center}

Como no exemplo acima, somos obrigados a utilizar parenteses sempre que for necessário evitar ambiguidades.
Se não os utilizarmos e simplesmente escrevermos ``$x\in y\text{ e } x=y \text{ ou } x\in z$'', não saberemos se a fórmula é $\left(x\in y\wedge x=y\right) \vee (x\in z)$ ou $x\in y\wedge \left(x=y \vee x\in z\right)$, que tem, intuitivamente, significados diferentes (reflita sobre a razão).

Novamente, vamos ilustrar o conceito de disjunção com um exemplo ``cotidiano''.

\begin{example}
    Considere a afirmação ``coelhos nascem de ovos'' e a afirmação ``vacinas previnem doenças''.

    Sua disjunção é a afirmação ``coelhos nascem de ovos ou vacinas previnem doenças''.

    Tal afirmação é verdadeira, uma vez que ao menos uma das afirmações que compõe a disjunção é verdadeira.
\end{example}

\subsection{O quantificador universal}
Se $\phi$ é uma fórmula e $x$ é uma variável lógica, então $\forall x \phi$ é a fórmula que lê-se ``para todo $x$, $\phi$'' ou ``para cada $x$, $\phi$''.

Ela é verdadeira se, para qualquer que seja o objeto $x$, a fórmula $\phi$ é verdadeira.
\begin{example}
    A fórmula $\forall x (x\notin \emptyset)$ é verdadeira, uma vez que para qualquer que seja o objeto $x$, $x$ não é um elemento do conjunto vazio $\emptyset$.
\end{example}

Novamente, vamos dar um exemplo ``cotidiano'' para ilustrar o conceito de quantificador universal.

\begin{example}
    A sentença ``todo mamífero é terrestre'' é falsa, uma vez que existem mamíferos aquáticos, como as baleias e golfinhos.

    Por outro lado, a sentença ``todo mamífero é um animal'' é verdadeira, uma vez que qualquer que seja o mamífero, ele é um animal por definição.
\end{example}

\subsection{O quantificador existencial}
Se $\phi$ é uma fórmula e $x$ é uma variável lógica, então $\exists x \phi$ é a fórmula que lê-se ``existe um $x$ tal que $\phi$'' ou ``existe pelo menos um $x$ tal que $\phi$''.

Ela é verdadeira se existe pelo menos um objeto $x$ tal que a fórmula $\phi$ é verdadeira para ele.

\begin{example}
    A fórmula $\exists x (x\in \emptyset)$ é falsa, uma vez que não existe nenhum objeto $x$ tal que $x$ seja um elemento do conjunto vazio $\emptyset$.

    Por outro lado, a fórmula $\exists x (x\in \mathbb R)$ é verdadeira, uma vez que existem objetos $x$ tais que $x$ é um elemento do conjunto dos números reais $\mathbb R$, como o número real $0$, por exemplo.
\end{example}

\begin{example}
    A sentença ``existe um mamífero que é aquático'' é verdadeira, uma vez que existem mamíferos aquáticos, como as baleias e golfinhos.

    Por outro lado, a sentença ``todo mamífero é aquático'' é falsa, uma vez que existem mamíferos terrestres, como os cães e gatos.
\end{example}
\subsection{O operador condicional}

Se $\phi$ e $\psi$ são fórmulas, então $\phi\to \psi$ é a fórmula condicional de $\phi$ para $\psi$.
Lemos $\phi\to \psi$ como ``se $\phi$, então $\psi$''.
Também podemos escrever $\phi\to \psi$ como literalmente ``se $\phi$, então $\psi$''.

A fórmula $\phi\to \psi$ é falsa quando $\phi$ é verdadeira e $\psi$ é falsa.
Nos outros casos, é verdadeira.

A tabela abaixo resume a relação entre os valores de verdade de $\phi$, $\psi$ e $\phi\to \psi$.

\begin{center}
    \begin{tabular}{cc|c}
        $\phi$ & $\psi$ & $\phi\to \psi$ \\
        \hline
        V & V & V \\
        V & F & F \\
        F & V & V \\
        F & F & V
    \end{tabular}
\end{center}

Tal convenção pode soar contra-intuitiva à primeira vista.
Para entender a razão por trás dessa convenção, convidamos o leitor a examinar o seguinte exemplo.
\begin{example}
    Considere a expressão ``para todo $x$, se $x$ é um número par, então $x$ é um número inteiro''.

    Intuitivamente, tal expressão é verdadeira: números pares são inteiros por definição.
    Porém, pelo exposto, para que $\forall x (x \text{ é um número par} \to x \text{ é inteiro})$ seja verdadeira, é necessário que para qualquer que seja o objeto $x$, a fórmula $x \text{ é um número par} \to x \text{ é inteiro}$ seja verdadeira.
    Ora, existem números que não são inteiros nem pares, como $\pi$.
    Para tais números, ambas as fórmulas $x \text{ é um número par}$ e $x \text{ é inteiro}$ são falsas.
    Apesar disso, a implicação $x \text{ é um número par} \to x \text{ é inteiro}$ é verdadeira para tais $x$.

    Notemos que se a tabela de verdade do condicional fosse diferente, a fórmula $\forall x (x \text{ é um número par} \to x \text{ é inteiro})$ seria falsa.
    A intuição é que os únicos $x$ que tornam uma implicação falsa são aqueles que servem como um contra-exemplo para a afirmação expressa pela implicação.
\end{example}

Notemos que se é verdadeira uma implicação da forma $\phi\to \psi$, e $\phi$ é verdadeira, então $\psi$ é verdadeira (do contrário, a implicação seria falsa).
Assim, de $\phi\to \psi$ e $\phi$, \emph{deduz-se} $\psi$.
Tal fenômeno é conhecido como \emph{modus ponens}, e é um dos princípios mais fundamentais da lógica.
Perceba que, se temos que $\phi\to \psi$ é verdadeira, e $\phi$ é falsa, então nada podemos concluir sobre $\psi$, uma vez que independentemente do valor de verdade de $\psi$, a implicação $\phi\to \psi$ é verdadeira, pela tabela.

\subsection{O operador bicondicional}
Se $\phi$ e $\psi$ são fórmulas, então $\phi\leftrightarrow \psi$ é a fórmula bicondicional de $\phi$ para $\psi$.
Lemos $\phi\leftrightarrow \psi$ como ``$\phi$ se e somente se $\psi$''.
Também podemos escrever $\phi\leftrightarrow \psi$ como literalmente ``$\phi$ se e somente se $\psi$'', ou, ainda, ``$\phi$ é equivalente a $\psi$''.

A fórmula $\phi\leftrightarrow \psi$ é verdadeira quando $\phi$ e $\psi$ têm o mesmo valor de verdade, ou seja, quando ambas são verdadeiras ou ambas são falsas.
Nos outros casos, é falsa.
A tabela abaixo resume a relação entre os valores de verdade de $\phi$, $\psi$ e $\phi\leftrightarrow \psi$.

\begin{center}
    \begin{tabular}{cc|c}
        $\phi$ & $\psi$ & $\phi\leftrightarrow \psi$ \\
        \hline
        V & V & V \\
        V & F & F \\
        F & V & F \\
        F & F & V
    \end{tabular}
\end{center}
\begin{example}
    A expressão ``existem unicórnios se, e somente se existem dragões'' é verdadeira, uma vez que ambas as afirmações ``existem unicórnios'' e ``existem dragões'' são falsas.
\end{example}

\subsection{Elementos da linguagem de conjuntos}
Dois conjuntos são iguais se possuem exatamente os mesmos elementos.
Escrevemos $x\subseteq y$ ($x$ é um subconjunto de $y$) para dizer que todo elemento de $x$ é um elemento de $y$.
Então $x=y$ se, e somente se, $x\subseteq y$ e $y\subseteq x$.

Se $A$ é um conjunto de elemetos $x$ e $\phi(x, A, t_1, \dots, t_n)$ é uma fórmula (que também pode depender de $t_1, \dots, t_n$), então $\{x\in A: \phi(x, A, t_1, \dots, t_n)\}$ é o conjunto de elementos $x$ de $A$ tais que $\phi(x, A, t_1, \dots, t_n)$ é verdadeira.
Por exemplo, $\{x\in \mathbb R: x>0\}$ é o conjunto de números reais $x$ tais que $x$ é maior do que $0$, ou seja, o conjunto dos números reais positivos.

Se $x_1, \dots, x_n$ são conjuntos, podemos formar o conjunto $\{x_1, \dots, x_n\}$, cujos elementos são exatamente $x_1, \dots, x_n$.
Note que, desse modo, $\{x_1, x_2\}=\{x_2, x_1\}$.

Se $\mathcal A$ é uma coleção de conjuntos, a união de $\mathcal A$ é o conjunto $\bigcup \mathcal A=\{x: \text{ existe } A\in \mathcal A \text{ tal que } x\in A\}$.
Ou seja, é o conjunto de todos os elementos de elementos de $A$.

Para visualizar tal conjunto, imagine que $\mathcal A$ é um carrinho de supermercado cheio de sacolas com produtos.
Cada sacola com produtos é um elemento de $\mathcal A$.
O conjunto $\bigcup \mathcal A$ é o conjunto de todos os produtos que estão dentro das sacolas do carrinho.
Para obtê-lo, basta abrir cada sacola e pegar os produtos que estão dentro dela, reunindo tudo em um novo conjunto (em uma caixa, por exemplo).

Se $A$ e $B$ são conjuntos, $A\cup B=\bigcup \{A, B\}$ é o conjunto de elementos de $A$ ou de $B$.
\[
    A\cup B=\{x: x\in A \text{ ou } x\in B\}.
\]
Lembremos que, em matemática, a palavra ``ou'' é inclusiva, ou seja, se $x$ é um elemento comum de $A$ e de $B$, então ainda assim $x$ é um elemento de $A\cup B$.

\begin{example}
    Se $A=\{1, 2, 3\}$ e $B=\{3, 4, 5\}$, então $A\cup B=\{1, 2, 3, 4, 5\}$.
\end{example}

Se $\mathcal A$ é uma coleção não vazia de conjuntos, a interseção de $\mathcal A$ é o conjunto $\bigcap \mathcal A=\{x: \text{ para todo } A\in \mathcal A, x\in A\}$.
Ou seja, é o conjunto de todos os elementos comuns a todos os elementos de $\mathcal A$.

Se $A$ e $B$ são conjuntos, $A\cap B=\bigcap \{A, B\}$ é o conjunto de elementos de $A$ e $B$ (simultaneamente).
\[
    A\cap B=\{x: x\in A \text{ e } x\in B\}.
\]

\begin{example}
    Se $A=\{1, 2, 3\}$ e $B=\{3, 4, 5\}$, então $A\cap B=\{3\}$.
\end{example}

Por fim, temos a diferença de conjuntos.
Se $A$ e $B$ são conjuntos, $A\setminus B=\{x: x\in A \text{ e } x\notin B\}$ é o conjunto de elementos de $A$ que não são elementos de $B$.
\begin{example}
    Se $A=\{1, 2, 3\}$ e $B=\{3, 4, 5\}$, então $A\setminus B=\{1, 2\}$.
\end{example}

Há um tipo especial de conjuntos chamado de \emph{par ordenado}.
Se $x$ e $y$ são pares ordenados, então $\langle x, y\rangle$ é o par ordenado de $x$ e $y$.
Formalmente, $x$ e $y$ não são, em geral, elementos de $\langle x, y\rangle$, de modo que não podemos escrever $x\in \langle x, y\rangle$ ou $y\in \langle x, y\rangle$.
Há construções formais específicas para criar tais objetos que não serão tratadas aqui.

Pares ordenados tem a propriedade de que $\langle x_1, y_1\rangle = \langle x_2, y_2\rangle$ se, e somente se, $x_1=x_2$ e $y_1=y_2$.
Dessa forma, diferente do que ocorre com $\{x_1, x_2\}$, a ordem dos termos que ocorrem no par ordenado é relevante.

\begin{example}
    $\langle 1, 2\rangle$ é um par ordenado de $1$ e $2$, e é diferente do par ordenado $\langle 2, 1\rangle$ de $2$ e $1$.

    Temos ainda que $\{1, 2\}\neq \langle 1, 2\rangle$.
    Enquanto o primeiro é o conjunto cujos únicos elementos são $1$ e $2$, o segundo possui outra constituição (não discutida aqui) e tem a propriedade fundamental de que, ao compará-lo com outro par ordenado, a ordem dos termos é relevante.
\end{example}

Se $A$ e $B$ são conjuntos, o produto cartesiano de $A$ e $B$ é o conjunto $A\times B=\{\langle x, y\rangle: x\in A \text{ e } y\in B\}$.
Uma relação entre $A$ e $B$ é um subconjunto de $A\times B$, ou seja, um conjunto de pares ordenados $\langle x, y\rangle$ tal que $x\in A$ e $y\in B$.
Uma função de $A$ em $B$ é uma relação $f\subseteq A\times B$ tal que, para cada $x\in A$, existe um único $y\in B$ tal que $\langle x, y\rangle \in f$.
Tal único $y$ é chamado de imagem de $x$ por $f$, e é denotado por $f(x)$.
Escrevemos $f: A\to B$ para dizer que $f$ é uma função de $A$ em $B$.
Neste contexto, $A$ é chamado de domínio de $f$, e $B$ é chamado de um contradomínio de $f$.
A imagem de $f$ é o conjunto $\{f(x): x\in A\}$, ou seja, o conjunto de todos os elementos de $B$ que são imagens de elementos de $A$ por $f$.
Note que todo contradomínio de $f$ contém a imagem de $f$, mas pode conter outros elementos que não são imagens de nenhum elemento de $A$ por $f$.

Uma função $f: A\to B$ é injetora se, para quaisquer $x_1, x_2\in A$, $f(x_1)=f(x_2)$ implica que $x_1=x_2$.
Uma função $f: A\to B$ é sobrejetora se, para cada $y\in B$, existe um $x\in A$ tal que $f(x)=y$.
Uma função $f: A\to B$ é bijetora se é injetora e sobrejetora.

Se $A, B, C$ são conjuntos e $f: A\to B$ e $g: B\to C$ são funções, então a composição de $f$ e $g$ é a função $g\circ f: A\to C$ tal que, para cada $x\in A$, $(g\circ f)(x)=g(f(x))$.

A função identidade de um conjunto $A$ é a função $\text{id}_A: A\to A$ tal que, para cada $x\in A$, $\text{id}_A(x)=x$.

Se $R\subseteq A\times B$ é uma relação, então a inversa de $R$ é a relação $R^{-1}=\{\langle y, x\rangle: \langle x, y\rangle \in R\}$.
Como toda função é uma relação, faz sentido considerar a inversa $f^{-1}$ de uma função $f$.
Fica como exercício para o leitor verificar que, $f^{-1}$ é uma função se, e somente se, $f$ é injetora.
Temos ainda que se $f: A\to B$ é uma função (com esse domínio e contradomínio), então $f^{-1}: B\to A$ é uma função (de $B$ em $A$), e somente se, $f$ é bijetora (de $A$ em $B$).

Se $f:A\rightarrow B$, a restrição de $f$ à $C$, onde $C\subseteq A$, é a função $f\restriction C: A\cap C\to B$ tal que, para cada $x\in A\cap C$, $f\restriction C(x)=f(x)$.
É muito comum fazer uso dessa notação quando $C$ já é, na verdade, um subconjunto do domínio de $f$, ou seja, quando $C\subseteq A$.
Nesses casos, $A\cap C=C$, e, portanto, $f\restriction C: C\to B$ é a função tal que, para cada $x\in C$, $f\restriction C(x)=f(x)$.

Se $A$ é um conjunto, uma operação binária em $A$ é uma função $f: A\times A\to A$.
Um exemplo usual é a soma de números reais, $+: \mathbb R\times \mathbb R\to \mathbb R$.
Em vez de escrevermos $+(\langle x, y\rangle)$, escrevemos $+(x, y)$, ou, muito mais normalmente, $x+y$.